\documentclass[11pt]{ctexart}
\usepackage[a4paper,margin=1in]{geometry}
\usepackage{graphicx}
\usepackage{xcolor}
\usepackage{hyperref}
\definecolor{refgray}{gray}{0.45}
\title{\textbf{市场最新观点汇总}\\[0.4em]{\large AI与能源双周期交织，因子选择优于区域押注，结构性分化主导资产定价}}
\author{KC桌面 自动生成}
\date{260522}
\begin{document}
\maketitle
\section*{一页摘要}
\begin{itemize}
  \item 全球市场正从贝塔驱动的同涨同跌转向因子驱动的结构性分化，AI资本开支与能源波动共同重塑盈利质量和增长可持续性的定价逻辑。
  \item 半导体设备支出在2028年逼近2000亿美元，核心驱动力从先进逻辑制程切换为存储芯片产能扩张，尤其是DRAM的结构性重投入。
  \item 黄金短期承压于能源通胀引发的实际利率和美元走强，黄金与实际利率的负相关性已恢复至-0.75，定价锚正在位移。
  \item 中国车企在欧洲的渗透从边缘试探进入核心市场突破阶段，英国市场单月销量同比增117\%，市占率升至15\%。
  \item 中国四月经济数据全面低于预期，但能源冲击和统计噪音解释了大部分下滑，政策空间的重新定价才是关键变量。
  \item 新兴市场AI采纳速度被低估，采纳滞后从五年缩至三年，开源模型和数字技术跳跃式发展正在压缩追赶周期。
  \item 中国天然气去库存速度超预期，储气量同比转负，下半年LNG进口压力将推升TTF和JKM的上行风险。
  \item 网络安全行业正经历AI驱动的供给侧结构性断裂，攻击成本下降500-1000倍，防御范式面临代际断层。
\end{itemize}
\section{宏观与利率}
\textbf{全球宏观环境正从总量增长叙事转向结构分化与政策空间重定价，能源冲击和AI投资周期交织，因子选择优先级高于区域选择。}

\begin{itemize}
  \item 美国实际GDP增速预计2.2\%，通胀从3.5\%逐步回落，美联储维持利率不变至2027年上半年，但企业间盈利能力的差距在拉大。
  \item 欧洲经济周期判断已从扩张调整为放缓，能源价格传导更快，央行反应函数不同，因子排序需转向现金流韧性和资本纪律。
  \item 中国四月经济数据综合偏离程度创2023年5月以来之最，但能源供给冲击可解释工业增加值约1个百分点的下滑，另有统计噪音因素。
  \item 中国政策制定者对经济状态的评估已发生变化，市场真正需要关注的是政策空间还有多大以及何时会被重新激活。
  \item 全球资本成本上升正在从边缘地带撕开裂口，新兴市场的信用裂缝可能通过汇率和通胀路径传导回核心资产定价。
\end{itemize}
\begin{figure}[htbp]
\centering
\includegraphics[width=0.88\linewidth]{figures/fig_144.jpg}
\caption{Exhibit 1 - Exhibit 1: April activity data was the largest downside surprise since May 2023}
\end{figure}
\begin{figure}[htbp]
\centering
\includegraphics[width=0.88\linewidth]{figures/fig_145.jpg}
\caption{Exhibit 2 - Exhibit 2: April physical output data shows the negative impact of the Iran war as well as continued strength in high-tech manufacturing and new energy sectors Bottom 10 sequential output growth from Mar to Apr}
\end{figure}
\begin{figure}[htbp]
\centering
\includegraphics[width=0.88\linewidth]{figures/fig_146.jpg}
\caption{Exhibit 3 - Exhibit 3: IP was strong in March, June and September but weak in April, July and October in 2025 China industrial production sequential growth (mom sa non-annualized)}
\end{figure}
\begin{figure}[htbp]
\centering
\includegraphics[width=0.88\linewidth]{figures/fig_151.jpg}
\caption{Exhibit 8 - Exhibit 8: April activity data continue to show a bifurcated economy with strong exports and weak domestic demand China April Activity Indicators}
\end{figure}
\begin{figure}[htbp]
\centering
\includegraphics[width=0.88\linewidth]{figures/fig_152.jpg}
\caption{Exhibit 9 - Exhibit 9: We expect sequential real GDP growth to slow from \$5.3\textbackslash{}\%\$ qoq annualized in Q1 to \$4.0\textbackslash{}\%\$ in Q2 China Current Activity Indicator and GDP}
\end{figure}
{\small\color{refgray}\textbf{References:} [R001] 市场真正低估的不是AI需求，而是能源波动如何重塑因子定价; [R005] 四月数据“惊吓”背后，市场真正需要关注的不是疲弱，而是政策空间的重新定价; [R013] 市场真正低估的不是AI泡沫，而是新兴市场的信用裂缝}

\section{AI与半导体}
\textbf{AI资本开支以远超预期的速度扩张，但半导体设备市场的核心驱动力已从先进逻辑制程切换为存储芯片产能扩张，硅电容等新市场正在起量。}

\begin{itemize}
  \item 全球WFE支出将在2028年逼近2000亿美元，DRAM是上修核心，2027年DRAM WFE从480亿美元上调至570亿美元，增幅近20\%。
  \item 中国存储扩产提速，长江存储和长鑫存储产能扩张信号积极，长鑫IPO后资金更充裕，将支持本地HBM制造。
  \item 硅电容市场正从AI服务器专用向高性能计算全场景扩张，一份约15亿美元的两年供应合同确认了应用边界的延伸。
  \item 硅电容制造不需要先进制程，传统DRAM产线即可生产，台湾DRAM厂有望受益于这一增量市场。
  \item AI对药物发现和CRO业务的影响已从概念走向可量化商业落地，金斯瑞生命科学业务已从AIDD中获得约50\%的新增蛋白合成需求。
\end{itemize}
\begin{figure}[htbp]
\centering
\includegraphics[width=0.88\linewidth]{figures/fig_053.jpg}
\caption{Exhibit 8 - EXHIBIT 1: We revise up our WFE projection for CY26-CY28 to reflect our latest views...}
\end{figure}
\begin{figure}[htbp]
\centering
\includegraphics[width=0.88\linewidth]{figures/fig_054.jpg}
\caption{EXHIBIT 3 - EXHIBIT 3: In China, we now see +15\% / +15\% YoY (+10\% / +7\% prior) in 2026 / 2027.}
\end{figure}
\begin{figure}[htbp]
\centering
\includegraphics[width=0.88\linewidth]{figures/fig_055.jpg}
\caption{EXHIBIT 2 - EXHIBIT 2: ...and we now model +21\% YoY in CY26 (+18\% prior). For CY27 we now model +18\% (+12\% prior) growth.}
\end{figure}
\begin{figure}[htbp]
\centering
\includegraphics[width=0.88\linewidth]{figures/fig_057.jpg}
\caption{EXHIBIT 5 - EXHIBIT 6: We now estimate 2027/2028 WFE spending for DRAM to be \textbackslash{}\$57bn / \textbackslash{}\$71bn (\textbackslash{}\$48bn / \textbackslash{}\$51bn prior) / NAND \textbackslash{}\$18bn / \textbackslash{}\$23bn (\textbackslash{}\$16bn / \textbackslash{}\$16bn prior), respectively...}
\end{figure}
\begin{figure}[htbp]
\centering
\includegraphics[width=0.88\linewidth]{figures/fig_058.jpg}
\caption{EXHIBIT 7 - EXHIBIT 7: ...and for Logic/Foundry, we keep our 2026-27 estimates, and lift 2028 to \textbackslash{}\$93bn (\textbackslash{}\$92bn prior).}
\end{figure}
{\small\color{refgray}\textbf{References:} [R002] 半导体设备市场的真正转折：2000亿美元WFE在望，但驱动力已彻底改变; [R011] 硅电容合同背后，真正被低估的是存储器工艺的重新定价; [R015] CXO 的真正拐点不是订单回暖，而是商业模式的重构}

\section{能源与大宗}
\textbf{能源波动成为宏观噪音的核心来源，天然气和黄金市场均面临供给侧再定价压力，中国补库需求将推升TTF和JKM上行风险。}

\begin{itemize}
  \item 中国4月储气量同比转负，为近年来首次，尽管需求疲弱，但LNG进口、国产气和管道气均低于预期，夏季补库压力将推高进口。
  \item 亚洲LNG价格（JKM）相对欧洲（TTF）溢价持续走强，吸引更多美国LNG船货从大西洋转向太平洋，西北欧LNG进口同比下滑。
  \item 如果霍尔木兹海峡流量正常化推迟至7月下旬，TTF价格可能分别上行至65和53欧元/兆瓦时，远超当前44-40欧元/兆瓦时的预测区间。
  \item 黄金与实际利率的负相关性已恢复至-0.75，能源价格推高实际利率和美元走强对黄金形成短期压力，0-3个月目标价4300美元/盎司。
  \item 中国黄金进口量仍维持在历史高位附近，年化约3000亿美元，一季度表观需求按价值计算创历史新高，但短期定价锚已位移。
\end{itemize}
\begin{figure}[htbp]
\centering
\includegraphics[width=0.88\linewidth]{figures/fig_160.jpg}
\caption{Exhibit 1 - Exhibit 1: China gas demand growth has remained weak... China apparent gas demand, Bcm/y}
\end{figure}
\begin{figure}[htbp]
\centering
\includegraphics[width=0.88\linewidth]{figures/fig_161.jpg}
\caption{Exhibit 2 - Exhibit 2: ...but not enough to prevent inventories moving down yoy in April China underground gas storage, Bcm}
\end{figure}
\begin{figure}[htbp]
\centering
\includegraphics[width=0.88\linewidth]{figures/fig_162.jpg}
\caption{Exhibit 3 - Exhibit 3: China LNG imports are increasing sequentially, which we expect to continue through the summer China LNG imports, mpta}
\end{figure}
\begin{figure}[htbp]
\centering
\includegraphics[width=0.88\linewidth]{figures/fig_163.jpg}
\caption{Exhibit 4 - Exhibit 4: The JKM-TTF premium remains strong enough to cover the incremental cost for the US to ship LNG to Asia vs to Europe JKM-TTF spread and the US shipping cost differential (lhs) vs US LNG loadings to Europe vs to Asia, mtpa}
\end{figure}
\begin{figure}[htbp]
\centering
\includegraphics[width=0.88\linewidth]{figures/fig_164.jpg}
\caption{Exhibit 5 - Exhibit 5: LNG loadings are now marginally up yoy as suppliers postpone maintenance... Global LNG loadings, mtpa}
\end{figure}
\begin{figure}[htbp]
\centering
\includegraphics[width=0.88\linewidth]{figures/fig_070.jpg}
\caption{Figure 1 - Alexander Hacking, CFA Figure 1. The negative correlation between gold price and US real rate has returned - average daily correlation between gold \& 5y real rate at -0.75 YTD, vs -0.65 in 2025 and only -0.2 in 2024}
\end{figure}
\begin{figure}[htbp]
\centering
\includegraphics[width=0.88\linewidth]{figures/fig_071.jpg}
\caption{Figure 2 - © 2026 Citi Inc. No redistribution without Citi's written permission. Figure 2. A major equity market correction typically sees gold lower first then rebound}
\end{figure}
\begin{figure}[htbp]
\centering
\includegraphics[width=0.88\linewidth]{figures/fig_075.jpg}
\caption{Figure 6 - © 2026 Citi Inc. No redistribution without Citi's written permission. Figure 6. China monthly non-monetary gold import set to remain strong on CNY strength, resilient domestic demand and easing import licensing rules}
\end{figure}
{\small\color{refgray}\textbf{References:} [R003] 黄金市场真正的短期压力不在需求，而在能源驱动的宏观逻辑切换; [R007] 中国补库压力被市场低估，TTF和JKM的上行风险正在积聚}

\section{地产与消费}
\textbf{中国消费和地产数据受政策节奏和短期扰动影响，结构分化明显，以旧换新补贴透支效应和能源价格冲击是主要拖累。}

\begin{itemize}
  \item 四月社零同比仅0.2\%，为2022年以来最低，主要原因是政府以旧换新补贴拉动的集中替换需求进入还债期，家电销售从+53\%降至-15\%。
  \item 新能源车购置税从0\%恢复到5\%，导致汽车销量明显下滑，但EV销售因汽油价格上涨而增加，ICE车销售下降。
  \item 工业增加值中化学制品和石油加工品产量环比下降最明显，与中东局势和能源价格走高直接相关，而工业机器人和风电设备产量增长显著。
  \item 中国PC出货量同比下滑2\%，平板出货量同比下滑3\%，但品牌间增速差接近100个百分点，华硕PC同比增38\%，华为PC同比暴跌47\%。
  \item 平板市场中苹果同比下滑21\%，联想平板逆势增长6\%，总量接近零增长下竞争从做大蛋糕变为切分蛋糕。
\end{itemize}
\begin{figure}[htbp]
\centering
\includegraphics[width=0.88\linewidth]{figures/fig_147.jpg}
\caption{Exhibit 4 - Exhibit 4: Effect of government-subsidized consumer goods trade-in program should fade over time}
\end{figure}
\begin{figure}[htbp]
\centering
\includegraphics[width=0.88\linewidth]{figures/fig_148.jpg}
\caption{Exhibit 5 - Exhibit 5: EV sales increased while ICE car sales declined in April on higher gasoline prices}
\end{figure}
\begin{figure}[htbp]
\centering
\includegraphics[width=0.88\linewidth]{figures/fig_149.jpg}
\caption{Exhibit 6 - Exhibit 6: FAI data diverged significantly from alternative measures of construction activity Real FAI Implied vs. Actual Commodity Demand}
\end{figure}
\begin{figure}[htbp]
\centering
\includegraphics[width=0.88\linewidth]{figures/fig_166.jpg}
\caption{Exhibit 1 - Exhibit 1: China PC shipments were down 2\% y/y to 2.56mn units in Mar-26}
\end{figure}
{\small\color{refgray}\textbf{References:} [R005] 四月数据“惊吓”背后，市场真正需要关注的不是疲弱，而是政策空间的重新定价; [R012] 中国PC与平板出货量下滑，真正值得关注的不是总量，而是份额的剧烈重塑}

\section{区域市场}
\textbf{中国资产重定价的核心是微观龙头的全球议价权，而非宏观反转；中国车企在欧洲和新兴市场AI采纳均出现结构性突破。}

\begin{itemize}
  \item 中国车企在欧洲3月单月销量达189,750辆，同比增长60\%，英国市场贡献超5.7万辆，同比增117\%，市占率从6\%升至15\%。
  \item 奇瑞、SAIC和比亚迪是欧洲市场三强，合计占中国品牌销量近70\%，英国右舵市场的成功验证了可复制的扩张模型。
  \item 新兴市场AI采纳速度被低估，采纳滞后从五年缩至三年，微软数据显示新兴市场AI使用率22\%，仅落后发达市场15个百分点。
  \item 开源模型和数字技术跳跃式发展正在压缩追赶周期，新兴市场企业在云服务等新业务领域更愿意直接跳到前沿技术。
  \item 全球资金对中国资产态度从全面回避转向极度挑剔的聚焦，AI半导体、工业自动化和电力设备等供给侧优势领域最受关注。
\end{itemize}
\begin{figure}[htbp]
\centering
\includegraphics[width=0.88\linewidth]{figures/fig_079.jpg}
\caption{Figure 2 - Figure 2: Chinese OEMs European market volume by automaker \# By Country Figure 3: European market total volume trend}
\end{figure}
\begin{figure}[htbp]
\centering
\includegraphics[width=0.88\linewidth]{figures/fig_081.jpg}
\caption{Figure 5 - Figure 5: Chinese OEMs' aggregate European market volume trend}
\end{figure}
\begin{figure}[htbp]
\centering
\includegraphics[width=0.88\linewidth]{figures/fig_082.jpg}
\caption{Figure 6 - Figure 6: Aggregate market share of Chinese OEMs in European market}
\end{figure}
\begin{figure}[htbp]
\centering
\includegraphics[width=0.88\linewidth]{figures/fig_093.jpg}
\caption{Figure 17 - Figure 17: Chinese OEMs' aggregate volume trend in the UK}
\end{figure}
\begin{figure}[htbp]
\centering
\includegraphics[width=0.88\linewidth]{figures/fig_094.jpg}
\caption{Figure 18 - Figure 18: Aggregate market share of Chinese OEMs in UK}
\end{figure}
\begin{figure}[htbp]
\centering
\includegraphics[width=0.88\linewidth]{figures/fig_155.jpg}
\caption{Exhibit 1 - Exhibit 1: AI Adoption in Major EMs Is Only Slightly Lagging DMs Microsoft AI Diffusion Share}
\end{figure}
\begin{figure}[htbp]
\centering
\includegraphics[width=0.88\linewidth]{figures/fig_156.jpg}
\caption{Exhibit 2 - Exhibit 2: The Absolute and Relative Gap Between DM and EM Adoption Rates Has Been Shrinking, Including for Cloud Services}
\end{figure}
\begin{figure}[htbp]
\centering
\includegraphics[width=0.88\linewidth]{figures/fig_158.jpg}
\caption{Exhibit 4 - Exhibit 4: EM Countries Are Much More Likely Leap to the Technology Frontier with Digital Technologies Share of EM Firms Using Most Advanced Technology for Business Function, by Business Function Type}
\end{figure}
{\small\color{refgray}\textbf{References:} [R004] 中国车企在欧洲的真正突破，不在销量，而在结构; [R006] 市场真正低估的不是AI需求，而是新兴市场的采纳速度; [R016] 中国资产重定价的核心不是宏观反转，而是微观龙头的全球议价权}

\section{风险偏好与资金流}
\textbf{市场情绪指标已触发反向信号，集中度超越历史泡沫时期，资金流向显示科技股狂热与防御资产流入并存。}

\begin{itemize}
  \item 某投行牛熊指标升至8.0，触发2002年以来第17次逆向卖出信号，历史平均2-3个月全球股市回调2\%-3\%，最大回撤可达15\%-20\%。
  \item AI相关前十大公司占美股市值约48\%，超过1970年代漂亮50和1990年代TMT泡沫时期的集中度，仅低于1880年代铁路股泡沫。
  \item 科技股单周流入90亿美元创去年10月来最大，但欧洲和新兴市场已连续6周资金流出，黄金流出11亿美元，加密货币流出15亿美元。
  \item 美债反而流入108亿美元为9周最多，显示资金在科技狂热中同时寻求避险，市场情绪高度分化。
  \item 网络安全行业攻击成本从国家级门槛降至任何人可负担，AI自主攻击事件已发生，防御范式正在被技术演进本身淘汰。
\end{itemize}
\begin{figure}[htbp]
\centering
\includegraphics[width=0.88\linewidth]{figures/fig_071.jpg}
\caption{Figure 2 - © 2026 Citi Inc. No redistribution without Citi's written permission. Figure 2. A major equity market correction typically sees gold lower first then rebound}
\end{figure}
\begin{figure}[htbp]
\centering
\includegraphics[width=0.88\linewidth]{figures/fig_072.jpg}
\caption{Figure 3 - © 2026 Citi Inc. No redistribution without Citi's written permission. Figure 3. Precious metals (particularly gold) historically outperformed during stagflationary periods (e.g. the \$2\textasciicircum{}\{nd\}\$ oil crisis), with volatility created by recessions / Volker moment co}
\end{figure}
{\small\color{refgray}\textbf{References:} [R013] 市场真正低估的不是AI泡沫，而是新兴市场的信用裂缝; [R017] 市场真正低估的不是需求，而是防御侧的代际断层}

\section{化工与材料}
\textbf{中国化工品贸易正经历定价权结构性转移，环氧树脂从以量换价转向以价定量，甲醇和锂盐均出现出口定价权超越进口依赖的拐点信号。}

\begin{itemize}
  \item 环氧树脂4月出口量价齐升，出口均价2560美元/吨同比涨24.6\%，但前四月累计出口量同比降2.5\%，供给端主动收缩后定价权修复。
  \item 中国环氧树脂净出口红利期收窄，4月净出口3.29万吨同比降25\%，为2023年以来最低水平，低价出货格局正在改善。
  \item 甲醇4月进口量同比降29\%至56.3万吨，但进口价格环比跳升32\%至365美元/吨，出口量同比飙升218\%至17.3万吨，量缩价涨格局值得跟踪。
  \item 中国氢氧化锂首次从净出口国转为净进口国，前四月净进口10.2千吨，而2025年同期为净出口35.4千吨，全球定价锚点正在位移。
  \item 碳酸锂进口量价齐升，4月进口均价16586美元/吨同比涨74\%，进口价格涨幅远超国内现货，海外供给溢价正在系统性地重新建立。
\end{itemize}
\begin{figure}[htbp]
\centering
\includegraphics[width=0.88\linewidth]{figures/fig_001.jpg}
\caption{Exhibit 2 - Exhibit 2: Backtest performance of our recommended US factors in rising and flat inflation scenarios – Internal Growth and EPS Revisions Breadth remain resilient across regimes, while Low PEG benefits from flatter inflation environ}
\end{figure}
\begin{figure}[htbp]
\centering
\includegraphics[width=0.88\linewidth]{figures/fig_002.jpg}
\caption{Exhibit 3 - Exhibit 3: Internal Growth and Low PEG remain resilient in high oil volatility, while EPS Revisions improve as volatility normalizes Annualized Performance of our recommended US Factors in High and Normal Oil Volatility Regimes}
\end{figure}
{\small\color{refgray}\textbf{References:} [R008] 中国环氧树脂出口的真正信号：价格弹性正在取代规模弹性; [R009] 中国甲醇贸易的拐点信号：出口定价权正在超越进口依赖; [R010] 中国锂贸易正在经历一次定价权的结构性转移，而非周期波动}

\section{储能与新能源}
\textbf{全球储能需求结构性增量发生板块轮动，欧洲正形成独立于中国的需求强化回路，大容量电芯和锂价上涨重塑成本优势。}

\begin{itemize}
  \item SMM将2026年全球储能电芯出货预测上调至926GWh，同比增长54\%，欧洲是最大增量，预计出货153GWh，同比增长94\%。
  \item 欧洲峰谷价差从约80欧元/MWh扩大至100欧元/MWh，是中国市场价差的两倍，辅助服务回报也在改善，项目IRR不再完全依赖补贴。
  \item 大容量电芯（500+Ah）开始被下游接受，锂价上涨使其成本优势更加明显，这是结构性变化而非短期波动。
  \item 碳酸锂涨价对2026年电芯产量的短期冲击有限，因成本传导链条长，电芯厂到集成商再到项目方存在时间差。
  \item 全球储能市场正从中国产能驱动全球成本下降的单边逻辑，转向区域需求分化、技术规格升级和成本传导机制重构的多维博弈。
\end{itemize}
\begin{figure}[htbp]
\centering
\includegraphics[width=0.88\linewidth]{figures/fig_001.jpg}
\caption{Exhibit 2 - Exhibit 2: Backtest performance of our recommended US factors in rising and flat inflation scenarios – Internal Growth and EPS Revisions Breadth remain resilient across regimes, while Low PEG benefits from flatter inflation environ}
\end{figure}
\begin{figure}[htbp]
\centering
\includegraphics[width=0.88\linewidth]{figures/fig_002.jpg}
\caption{Exhibit 3 - Exhibit 3: Internal Growth and Low PEG remain resilient in high oil volatility, while EPS Revisions improve as volatility normalizes Annualized Performance of our recommended US Factors in High and Normal Oil Volatility Regimes}
\end{figure}
{\small\color{refgray}\textbf{References:} [R014] 全球储能市场真正超预期的不是中国，而是欧洲的自我强化循环}

\section*{结语}
综合来看，当前市场主线是AI与能源双周期交织下的结构性分化，因子选择优于区域押注，投资者需关注供给侧能力重定价、新兴市场追赶加速以及防御范式的代际断层。
\vfill
{\small\color{refgray}Personal reading notes and learning share only. Not investment advice.}
\end{document}
